\chapter{Fentanyl}

\section*{Definition and Overview}
Fentanyl is a powerful synthetic opioid medicine, roughly 100 times stronger than morphine, used for the relief of severe pain. In medical practice it is given as an injection, patch, lozenge, or nasal spray, and it is highly effective for pain during surgery, in intensive care, for cancer pain, and in palliative care. Because of its potency, fentanyl also carries a high risk of overdose, and in recent years illicitly manufactured fentanyl and its analogues have been mixed into heroin and other street drugs, causing a global surge in overdose deaths. Used exactly as prescribed, fentanyl is a safe and valuable medicine; used outside medical supervision, it is extremely dangerous. This chapter explains how fentanyl works, how it is used, and the risks that demand respect.

\section*{Epidemiology}
Fentanyl is one of the most commonly used opioids in Australian hospitals, and its legitimate medical use has grown over recent decades. Alongside this, illicit fentanyl has driven an international overdose crisis: in North America it is now the leading cause of drug overdose deaths, and it has been detected increasingly in Australian drug markets. Opioid overdoses kill more Australians than many other causes of injury, and fentanyl contributes to these deaths because of its potency and because users may not know it is present in the drugs they take. Fentanyl patches are also a common source of diversion and accidental exposure, including in children.

\section*{Aetiology and Pathophysiology}
\begin{itemize}
    \item \textbf{Mechanism:} fentanyl binds strongly to opioid receptors in the brain and spinal cord, blocking pain signals and producing analgesia, sedation, and euphoria
    \item \textbf{Potency:} because it is about 100 times more potent than morphine, effective doses are tiny and the margin between a therapeutic dose and a fatal dose is narrow
    \item \textbf{Rapid onset:} intravenous and smoked fentanyl act within seconds to minutes, and even transdermal absorption can be rapid in overdose
    \item \textbf{Respiratory depression:} the dangerous effect of opioids is suppression of the brain's drive to breathe; at overdose doses this leads to slowed breathing, hypoxia, and death
    \item \textbf{Tolerance and dependence:} regular use leads to tolerance (needing more for the same effect) and physical dependence, with withdrawal on stopping
    \item \textbf{Illicit use:} fentanyl is added to heroin, cocaine, counterfeit pills, and other drugs as a cheap way to increase potency, without users knowing
    \item \textbf{Route matters:} patches release fentanyl slowly over 72 hours, while lozenges and sprays are for breakthrough pain; misuse of patches, including chewing or applying heat, can be fatal
\end{itemize}

\section*{Clinical Features}
The effects and signs of overdose are specific.
\begin{itemize}
    \item \textbf{Intended effects:} profound pain relief, calm, and sedation
    \item \textbf{Common side effects:} drowsiness, nausea, vomiting, constipation, and confusion
    \item \textbf{Signs of overdose:} severe drowsiness or unresponsiveness, slow or shallow breathing, tiny (pinpoint) pupils, cold clammy skin, and limpness
    \item \textbf{Blue lips or skin:} a sign of dangerous oxygen lack
    \item \textbf{Snoring or gurgling breathing:} can indicate obstructed or failing breathing in overdose
    \item \textbf{Withdrawal on stopping:} agitation, muscle aches, sweating, runny nose, diarrhoea, and strong craving
    \item \textbf{Children:} accidental exposure to a patch can cause drowsiness and breathing difficulty
\end{itemize}

\section*{Differential Diagnosis}
Breathing problems and drowsiness in opioid exposure should be distinguished from other causes.
\begin{itemize}
    \item Overdose of other sedatives, including benzodiazepines and alcohol
    \item Other medical causes of reduced consciousness, including stroke, head injury, low blood sugar, and sepsis
    \item Respiratory disease causing hypoxia, such as pneumonia or COPD
    \item Opioid toxicity from other opioids such as morphine, oxycodone, or methadone
    \item Non-opioid drug effects causing agitation or confusion
\end{itemize}

\section*{When to Seek Medical Care}
Anyone prescribed fentanyl should be educated on correct use, storage, and disposal, and should seek medical review if pain is not controlled, side effects are troublesome, or there are signs of excessive sedation. Anyone who uses drugs should be aware of the risk of fentanyl contamination and should not use alone.

\textbf{Call 000 (or local emergency services) for:}
\begin{itemize}
    \item Signs of overdose: unresponsiveness, slow or stopped breathing, or tiny pupils
    \item Snoring, gurgling, or irregular breathing in a drowsy person
    \item A child who may have touched or swallowed a fentanyl patch
    \item Blue lips, cold clammy skin, or collapse
    \item Severe drowsiness with difficulty waking
\end{itemize}

\section*{Diagnosis and Investigations}
\begin{itemize}
    \item \textbf{Clinical recognition:} the combination of sedation, respiratory depression, and pinpoint pupils in a person with opioid exposure is the key to rapid diagnosis
    \item \textbf{Response to naloxone:} administration of the antidote naloxone is both diagnostic and life-saving; improvement confirms opioid toxicity
    \item \textbf{Toxicology testing:} urine or blood drug screens may confirm exposure, though fentanyl analogues can be missed by standard tests
    \item \textbf{Blood gases:} oxygen and carbon dioxide levels show the severity of respiratory failure
    \item \textbf{Investigations for complications:} chest X-ray for aspiration, and ECG for rhythm effects
\end{itemize}

\section*{Management}
\begin{itemize}
    \item \textbf{Medical use:} given by trained staff in hospital; patches and lozenges prescribed for chronic severe pain under specialist supervision, with strict instructions
    \item \textbf{Overdose emergency:} call 000 immediately; give naloxone if available (including take-home naloxone programs); place the person on their side; begin rescue breathing if trained; stay until help arrives because naloxone can wear off before fentanyl does
    \item \textbf{Hospital treatment:} airway support, oxygen, repeated naloxone, and observation for several hours because of fentanyl's long duration of action
    \item \textbf{Safe prescribing:} lowest effective dose, review of risks, and limiting quantities; fentanyl is avoided in people who misuse opioids
    \item \textbf{Dependence treatment:} opioid dependence is treated with methadone or buprenorphine and psychosocial support
    \item \textbf{Prevention of diversion:} secure storage, counted supplies, and safe disposal of used patches
    \item \textbf{Harm reduction:} drug checking services, take-home naloxone, and not using alone reduce deaths from contaminated drugs
\end{itemize}

\section*{Complications}
\begin{itemize}
    \item Fatal respiratory depression from overdose, including from accidental exposure
    \item Brain damage from prolonged hypoxia if breathing is not supported
    \item Aspiration pneumonia from vomiting while unconscious
    \item Dependence, addiction, and withdrawal
    \item Constipation, which can be severe with long-term use
    \item Interactions with alcohol and other sedatives, which increase overdose risk
    \item The risks of the street drug trade, including unknowingly taking fentanyl
\end{itemize}

\section*{Prognosis and Outcomes}
As a medical treatment, fentanyl is very effective, and with appropriate prescribing and monitoring, serious harm is uncommon. The outlook in overdose depends on speed of response: with rapid naloxone and breathing support, most people survive, but delayed help or use of fentanyl with other depressants is often fatal. For people who develop opioid dependence, treatment with methadone or buprenorphine greatly improves survival and quality of life. Public health measures, including naloxone access and drug checking, are reducing deaths from illicit fentanyl.

\section*{Prevention and Self-Care}
\begin{itemize}
    \item Use fentanyl only exactly as prescribed, and never share or sell it
    \item Never combine fentanyl with alcohol, benzodiazepines, or other sedatives
    \item Store patches and medicines securely out of reach of children and pets
    \item Dispose of used patches by folding them and returning them to the pharmacy
    \item If you use drugs, never use alone, and carry naloxone if you can obtain it through take-home programs
    \item Know the signs of overdose and how to call for help and give rescue breathing
    \item Seek help early for pain management or substance use concerns
\end{itemize}

\section*{Sources and Further Reading}
The following reputable sources provide further information for healthcare students and patients seeking reliable, current advice about fentanyl.
\begin{itemize}
    \item Healthdirect Australia. \textit{Fentanyl.} https://www.healthdirect.gov.au/
    \item Therapeutic Goods Administration. \textit{Fentanyl medicines information.} https://www.tga.gov.au/
    \item Penington Institute. \textit{Harm reduction and naloxone resources.} https://www.penington.org.au/
    \item Alcohol and Drug Foundation. \textit{Fentanyl facts.} https://adf.org.au/
    \item NHS. \textit{Fentanyl.} https://www.nhs.uk/medicines/
    \item United Nations Office on Drugs and Crime. \textit{Fentanyl and synthetic drugs.} https://www.unodc.org/
\end{itemize}
